\documentclass[10pt,a4paper]{amsart}
\usepackage[margin=19mm]{geometry}
\usepackage{amsmath,amssymb,mathtools}
\usepackage[scr=rsfs,cal=euler]{mathalfa}
\usepackage{xfrac}
\usepackage[unicode,hidelinks,bookmarks=false]{hyperref}
\newcommand{\ie}{i.e.\ }
\newcommand{\cf}{cf.\ }
\newcommand{\resp}{respectively\ }
\newcommand{\Subsection}{\subsection}
\providecommand{\nobreakdash}{\nobreak\hbox{-}}
\setlength{\emergencystretch}{2em}
\multlinegap0pt
\usepackage{amssymb}
\usepackage{csquotes}
\usepackage[shortlabels]{enumitem}
\newtheorem{lemma}{Lemma}
\title{The index of a numerical semigroup ring}
\author{Richard Bartels, Sarah Dajani,, Gabriel Koomson}
\date{Source: arXiv:2608.26003v2 (CC BY 4.0)}
\begin{document}
\maketitle

\noindent\emph{Source and licence.} The index of a numerical semigroup ring. Version arXiv:2608.26003v2 (CC BY 4.0), distributed by arXiv under the Creative Commons Attribution 4.0 International licence, http://creativecommons.org/licenses/by/4.0/, as recorded in the arXiv licence metadata for this exact version and shown on the abstract page https://arxiv.org/abs/2608.26003v2. Full licence notice: \texttt{licenses/arxiv-2608.26003-CC-BY-4.0.txt}.\par\smallskip

\noindent\textit{Editorial scope.} Counted unit: the article's lemma lemma (source0.tex lines 474--479). The article's complete original proof of it is delivered with it (source0.tex lines 480--491, one \texttt{proof} environment of 12 lines). No mathematics is written, repaired or re-derived: the statement and the proof are the article's own lines, in the article's own notation, and no retained line is substituted or re-flowed.  No further source-local context is needed: every cross-reference of every retained line resolves inside the two delivered spans.  Nothing was omitted from any retained span, so the package carries no omission record.  No retained line contains a citation, so the wrapper carries no bibliography and no bibliography entry.  No retained line contains a figure environment, an image-inclusion command or drawing code, so no image is delivered and the package declares no asset dependency.  Editorial formatting only, in addition: an AMS \texttt{amsart} preamble that restates the article's own declarations and macros used by the retained lines, namely the environments \texttt{lemma} on separate counters, together with the standard packages those lines need.  The retained environments print in this document as Lemma 1 in order of appearance, whereas the article prints the same items with the numbering of its own full document.  The complete native source of the article as distributed in the arXiv e-print for this exact version is bundled unchanged as \texttt{sources/208613/source0.tex}.
\begin{lemma}
Let $p$ be a prime integer. Define the polynomial $f(x,y) \in \mathbb{F}_{p}[x,y]$ by 
\[
f(x,y) = y \left(\prod\limits_{\alpha \in \mathbb{F}_p} (x+\alpha y) \right).
\] Then $f(x,y)=x^p y-xy^p$.
\end{lemma}\text{}\\

\begin{proof}
Let $g(x)=x^p-x \in \mathbb{F}_{p}[x]$. For every element $\alpha \in \mathbb{F}_{p}$, we have $g(\alpha)=0$. Therefore, $(x-\alpha) \mid g(x)$ for all $\alpha \in \mathbb{F}_{p}$. Since the linear polynomials $x-\alpha_1$ and $x-\alpha_2$ are coprime in $\mathbb{F}_p[x]$ for distinct elements $\alpha_1 \neq \alpha_2$ of $\mathbb{F}_p$, we have \\
\[
\left(\prod\limits_{\alpha \in \mathbb{F}_p} (x+\alpha ) \right) \Biggm| g(x)  
\] \text{}\\ and $g(x)=\prod\limits_{\alpha \in \mathbb{F}_p} (x+\alpha)$. Then \\\\
\[
x^p-xy^{p-1}=y^pg(x/y)=\prod\limits_{\alpha \in \mathbb{F}_p} (x+\alpha y) 
\] and 
\[
f(x,y)=x^p y- xy^p.
\]
\end{proof}

\end{document}
